% !TEX root = ../main.tex

\section{Analysis} \label{sec:analysis}
Recalling the reaction equation:
$$\rxnpl$$
The identity of the precipitate can be either \ce{BaSO4} or \ce{AlCl3}.

\ce{AlCl3} can be immediately proved to not be the precipitate. In the second trial, the following reaction occurred:
$$\ce{3BaCl2} + \ce{2Al(NO3)3} \rightleftharpoons \ce{3Ba(NO3)2} + \ce{2AlCl3}$$
Because no precipitate formed, we know that both \ce{Ba(NO3)2} and \ce{AlCl3} are soluble in water. Therefore, \ce{AlCl3} cannot be the precipitate.

\ce{BaSO4} can be proved to be the precipitate by observing the reactions and noting that a precipitate only forms when both \ce{Ba^2+} and \ce{SO4^2-} are present in the reagents.
In the first and third trials, both are present, and both form precipitates. In the second trial, \ce{SO4^2-} is absent, and no precipitate is formed. In the fourth trial, \ce{Ba^2+} is absent, and no precipitate is formed. The absence of either leads to no precipitate forming, meaning that both ions must be present in the percipitate. Combining them, the precipitate can be deduced to be \ce{BaSO4}.